\editionPage[\pdfbookmark[0]{Deux brevets blancs}{bm-bilans-guide-parent}]{bilans-guide}{Deux brevets blancs}{Session 2027 · Sujets originaux}{Durée totale : 2 heures · Note sur 20}
\begin{blocEdition}{bilans-guide}

L’épreuve de mathématiques du DNB, série générale, est affectée du coefficient 2.
\editionEtape{Partie 1 : 20 minutes, sans calculatrice}
Automatismes et applications directes : 6 points. Douze questions à réponse courte,
0,5 point chacune. Préparer une première copie, rendue après 20 minutes.
\editionEtape{Partie 2 : 1 h 40, calculatrice autorisée}
Raisonnement et résolution de problèmes : 14 points. Quatre problèmes
indépendants à rédiger, sans QCM. Utiliser une seconde copie. Le passage
de la première à la seconde partie se fait sans pause.
\editionEtape{Le barème de ce livre}
Chaque problème vaut 3,5 points : trois questions à 1 point et 0,5 point
pour la clarté, les unités et la rédaction. Les 2 points de qualité
d’expression sont ainsi inclus dans les 14 points de problèmes.
Ce partage détaillé est un choix d’entraînement ; il ne préjuge pas
des barèmes des futurs sujets officiels.
\editionEtape{Après la séance}
Comparer les démarches, classer les erreurs par notion et refaire les
questions utiles sans correction. Un résultat numérique faux n’empêche
pas de valoriser une démarche correcte.

\end{blocEdition}
\editionPage{bilan-A-auto}{Brevet blanc A}{Relier calculs, modèles et figures}{Partie 1 · 20 min · Sans calculatrice · 6 points}
\begin{blocEdition}{bilan-A-auto}
\begin{questions}
\item Calculer $-7+3\times5$.
\item Calculer $2/3+1/6$.
\item Donner l’écriture scientifique de 450000.
\item Donner $\sqrt{81}$.
\item Développer $3(x-4)$.
\item Résoudre $5x=35$.
\item Calculer 15 \% de 80.
\item Pour $f(x)=2x+1$, calculer $f(-2)$.
\item Convertir 0,25 h en minutes.
\item Donner la médiane de $2;4;5;8;11$.
\item On tire au hasard dans un sac de 2 rouges et 3 bleues. Donner la probabilité de bleu.
\item Un carré a un côté de 3 cm. Donner son aire.\end{questions}
\end{blocEdition}
\editionPage{bilan-A-ex-1}{Brevet A · Problème 1}{Un atelier de carnets}{Partie 2 · 1 h 40 pour les quatre problèmes · 3,5 points}
\begin{blocEdition}{bilan-A-ex-1}
Un atelier possède 168 feuilles bleues et 126 feuilles blanches. Il prépare le plus grand nombre possible de carnets identiques, sans reste. Chaque carnet contient le même nombre de feuilles de chaque couleur.\begin{questions}
\item Déterminer le nombre maximal de carnets et leur composition.
\item L’atelier vend chaque carnet 4,50 €. Les fournitures ont coûté 78 €. Calculer le solde si tous les carnets sont vendus.
\item Un autre atelier vend le carnet 5 €, puis applique 10 \% de réduction. Comparer son prix unitaire au précédent.
\end{questions}\par\medskip{\small\color{otMuted}Barème : question 1 : 1 pt ; question 2 : 1 pt ; question 3 : 1 pt ; clarté de la rédaction : 0,5 pt.}
\end{blocEdition}
\editionPage{bilan-A-ex-2}{Brevet A · Problème 2}{Une voile divisée}{Partie 2 · 1 h 40 pour les quatre problèmes · 3,5 points}
\begin{blocEdition}{bilan-A-ex-2}
ABC est rectangle en A, AB = 12 m, AC = 9 m. M appartient à [AB] avec AM = 4 m. N appartient à [AC] et MN est parallèle à BC.\begin{questions}
\item Calculer BC.
\item Calculer AN et MN.
\item Calculer l’aire de ABC puis celle du quadrilatère MBCN.
\end{questions}\par\medskip{\small\color{otMuted}Barème : question 1 : 1 pt ; question 2 : 1 pt ; question 3 : 1 pt ; clarté de la rédaction : 0,5 pt.}\label{fig-bilan-A-ex-2}\figureGeo{brevet-A-geo}

\end{blocEdition}
\editionPage{bilan-A-ex-3}{Brevet A · Problème 3}{Deux locations de matériel}{Partie 2 · 1 h 40 pour les quatre problèmes · 3,5 points}
\begin{blocEdition}{bilan-A-ex-3}
Pour n journées entières de location, A facture $18+7n$ euros et B facture $10n$ euros, pour $1\le n\le15$.\begin{questions}
\item Calculer les deux prix pour 4 jours. Quelle offre choisir ?
\item Pour combien de jours les prix sont-ils égaux ? Préciser l’offre la moins chère avant et après ce nombre.
\item Avec 90 €, déterminer la durée maximale possible chez chaque loueur. Justifier les deux maximums.
\end{questions}\par\medskip{\small\color{otMuted}Barème : question 1 : 1 pt ; question 2 : 1 pt ; question 3 : 1 pt ; clarté de la rédaction : 0,5 pt.}\label{fig-bilan-A-ex-3}\figureGeo{brevet-A-tarifs}

\end{blocEdition}
\editionPage{bilan-A-ex-4}{Brevet A · Problème 4}{Des essais et un tirage}{Partie 2 · 1 h 40 pour les quatre problèmes · 3,5 points}
\begin{blocEdition}{bilan-A-ex-4}
Une série de 20 essais donne les scores du tableau. On tire ensuite deux boules avec remise dans un sac de 3 rouges et 1 bleue. Les tirages sont réalisés au hasard, avec mélange à chaque étape.\begin{center}\renewcommand{\arraystretch}{1.35}\begin{tabular}{c|cccc}\toprule
Score & $0$ & $1$ & $2$ & $3$\\\midrule
Effectif & $2$ & $5$ & $8$ & $5$\\
\bottomrule\end{tabular}\end{center}\begin{questions}
\item Calculer la moyenne et la médiane des scores.
\item Calculer la probabilité d’obtenir deux rouges.
\item Calculer la probabilité d’obtenir exactement une bleue. Comparer à la probabilité de deux rouges.
\end{questions}\par\medskip{\small\color{otMuted}Barème : question 1 : 1 pt ; question 2 : 1 pt ; question 3 : 1 pt ; clarté de la rédaction : 0,5 pt.}\label{fig-bilan-A-ex-4}\figureGeo{brevet-A-arbre}

\end{blocEdition}
\editionPage{bilan-B-auto}{Brevet blanc B}{Choisir une méthode et contrôler}{Partie 1 · 20 min · Sans calculatrice · 6 points}
\begin{blocEdition}{bilan-B-auto}
\begin{questions}
\item Calculer $18-4\times3$.
\item Calculer $(3/5)\div(9/10)$.
\item Écrire $10^3\times10^{-5}$ comme une puissance de 10.
\item Décomposer 60 en facteurs premiers.
\item Factoriser $x^2-49$.
\item Résoudre $x+8=3$.
\item Augmenter 50 de 12 \%.
\item Pour $g(x)=-3x$, donner l’antécédent de 12.
\item Convertir 2,4 m$^3$ en litres.
\item Donner l’étendue de $-2;1;4;9$.
\item Un événement a une probabilité de 0,3. Donner celle de son contraire.
\item Une longueur de 8 cm est réduite avec un rapport de $-0{,}5$. Donner la longueur de son image.\end{questions}
\end{blocEdition}
\editionPage{bilan-B-ex-1}{Brevet B · Problème 1}{Un abri et sa pente}{Partie 2 · 1 h 40 pour les quatre problèmes · 3,5 points}
\begin{blocEdition}{bilan-B-ex-1}
Un toit est représenté par ABC rectangle en A. AB = 4 m horizontalement, AC = 2,5 m verticalement et BC est le versant. L’abri a une longueur de 6 m perpendiculaire au dessin.\begin{questions}
\item Calculer BC au centimètre.
\item Calculer l’angle B au degré.
\item Le versant rectangulaire doit être couvert. Ajouter 8 \% à son aire pour les découpes ; chaque plaque couvre 2 m$^2$. Combien de plaques au minimum ? Utiliser la longueur BC non arrondie.
\end{questions}\par\medskip{\small\color{otMuted}Barème : question 1 : 1 pt ; question 2 : 1 pt ; question 3 : 1 pt ; clarté de la rédaction : 0,5 pt.}\label{fig-bilan-B-ex-1}\figureGeo{brevet-B-geo}

\end{blocEdition}
\editionPage{bilan-B-ex-2}{Brevet B · Problème 2}{Une bouteille et un débit}{Partie 2 · 1 h 40 pour les quatre problèmes · 3,5 points}
\begin{blocEdition}{bilan-B-ex-2}
Une bouteille est modélisée par un cylindre de rayon intérieur 4 cm et de hauteur utile 20 cm. Une machine la remplit à 50 mL/s, à débit constant.\begin{questions}
\item Calculer la capacité exacte en cm$^3$, puis l’arrondir au mL.
\item Calculer la durée nécessaire pour remplir 90 \% de la capacité, au dixième de seconde.
\item Une version réduite a toutes ses longueurs multipliées par 0,75. Calculer sa capacité au mL et expliquer pourquoi elle ne vaut pas 75 \% de l’ancienne.
\end{questions}\par\medskip{\small\color{otMuted}Barème : question 1 : 1 pt ; question 2 : 1 pt ; question 3 : 1 pt ; clarté de la rédaction : 0,5 pt.}\label{fig-bilan-B-ex-2}\figureGeo{brevet-B-cylindre}

\end{blocEdition}
\editionPage{bilan-B-ex-3}{Brevet B · Problème 3}{Un script de figures}{Partie 2 · 1 h 40 pour les quatre problèmes · 3,5 points}
\begin{blocEdition}{bilan-B-ex-3}
Le projet démarre stylo baissé, orienté vers la droite. Le bloc carré(c) trace un carré de côté c et revient au départ avec l’orientation initiale. Le programme met c à 10, puis répète quatre fois : carré(c) ; relever le stylo ; avancer de c + 5 ; baisser le stylo ; ajouter 5 à c.\begin{questions}
\item Donner les côtés successifs des quatre carrés.
\item Calculer la somme des périmètres tracés.
\item Donner la valeur finale de c et la distance totale avancée stylo relevé. Expliquer pourquoi la valeur finale de c n’est pas le côté d’un cinquième carré tracé.
\end{questions}\par\medskip{\small\color{otMuted}Barème : question 1 : 1 pt ; question 2 : 1 pt ; question 3 : 1 pt ; clarté de la rédaction : 0,5 pt.}\label{fig-bilan-B-ex-3}\figureGeo{brevet-B-script}

\end{blocEdition}
\editionPage{bilan-B-ex-4}{Brevet B · Problème 4}{Un jeu à deux étapes}{Partie 2 · 1 h 40 pour les quatre problèmes · 3,5 points}
\begin{blocEdition}{bilan-B-ex-4}
Une roue possède quatre secteurs égaux : trois portent A, un porte B. Un sac contient une boule rouge et une bleue. On tourne la roue, puis on tire une boule au hasard, indépendamment de la roue. Le joueur gagne sur A-rouge ou B-bleue.\begin{questions}
\item Calculer la probabilité de chacun des deux chemins gagnants.
\item En déduire la probabilité de gagner. Sur 100 parties, peut-on garantir exactement 50 gains ?
\item Un gain rapporte 6 points, la participation coûte 2 points. Selon la règle « moyenne théorique = somme des bilans multipliés par leurs probabilités », calculer le bilan net moyen par partie.
\end{questions}\par\medskip{\small\color{otMuted}Barème : question 1 : 1 pt ; question 2 : 1 pt ; question 3 : 1 pt ; clarté de la rédaction : 0,5 pt.}\label{fig-bilan-B-ex-4}\figureGeo{brevet-B-arbre}

\end{blocEdition}
